\documentclass[11pt]{article}

\usepackage[T1]{fontenc}
\usepackage[utf8]{inputenc}
\usepackage{lmodern}
\usepackage[margin=1.04in]{geometry}
\usepackage{amsmath,amssymb,mathtools,amsthm}
\usepackage{mathrsfs}
\usepackage{microtype}
\usepackage{enumitem}
\usepackage{xcolor}
\usepackage{booktabs}
\usepackage{array}
\usepackage{longtable}
\usepackage{hyperref}
\usepackage[capitalize,noabbrev]{cleveref}

\hypersetup{
  colorlinks=true,
  linkcolor=blue!45!black,
  citecolor=green!35!black,
  urlcolor=blue!55!black,
  pdftitle={Structural Companion to Algebraic Exponential Integrals},
  pdfauthor={Christopher D. Long}
}

\numberwithin{equation}{section}
\setlength{\emergencystretch}{3em}

\newtheorem{theorem}{Theorem}[section]
\newtheorem{proposition}[theorem]{Proposition}
\newtheorem{lemma}[theorem]{Lemma}
\newtheorem{corollary}[theorem]{Corollary}
\theoremstyle{definition}
\newtheorem{definition}[theorem]{Definition}
\theoremstyle{remark}
\newtheorem{remark}[theorem]{Remark}

\newcommand{\C}{\mathbb C}
\newcommand{\Q}{\mathbb Q}
\newcommand{\Qbar}{\overline{\mathbb Q}}
\newcommand{\Aone}{\mathbb A^1}
\newcommand{\Span}{\operatorname{span}}
\newcommand{\Spec}{\operatorname{Spec}}
\newcommand{\ord}{\operatorname{ord}}
\newcommand{\Mcal}{\mathcal M}
\newcommand{\Ecal}{\mathcal E}
\newcommand{\Rcal}{\mathcal R}
\newcommand{\Dmod}{\mathcal D}
\newcommand{\Ocal}{\mathcal O}
\newcommand{\Lcal}{\mathscr L}
\newcommand{\Vcal}{\mathscr V}

\title{Structural Companion to Algebraic Exponential Integrals:\\
Audit, Detailed Lemmas, and Directions for Generalization}
\author{Christopher D. Long}
\date{August 2026}

\begin{document}
\maketitle

\begin{abstract}
This companion records the material removed from the streamlined base paper
\emph{Algebraic Exponential Integrals: Transcendence and Linear
Independence}.  It has three purposes.  First, it audits the earlier
1,832-line draft and records the parameter absorptions and proof compression.
Second, it preserves the stronger structural results that are not needed for
the shortest base proof: the full Brieskorn lattice, a detailed
Fourier--Laplace semisimplicity argument, the complete horizontal
endomorphism algebra, multiplicity-free decomposition, the relation
differential submodule of relations and its constant projector, and the exact
Cauchy-transform criterion.  Third, it separates polynomial-specific
calculations from mechanisms likely to survive for finite maps of curves,
larger families of exponential period functions, and algorithmic
$E$-function problems.
\end{abstract}

\tableofcontents

\section{Vocabulary and scope}

The revised base paper follows the terminology used in the closest
$E$-function and exponential-period literature.  A parameter integral
$\int_\gamma e^{-zf}\omega$ is an \emph{exponential period function}; $f$
is a regular function and, in the one-variable case considered here, a
polynomial \cite{FresanJossen,CommelinHabeggerHuber}.  The geometric language
is that of exponential twists, twisted de~Rham cohomology,
Fourier--Laplace transforms, and topological or relative chains
\cite{HienRoucairol,SabbahYu}.  For the final analytic step, the established
terms are moments, the polynomial moment problem, Cauchy transforms,
zero-cycles, inverse branches, and polynomial constellations
\cite{PakovichMuzychuk,GavrilovPakovich}.

Accordingly, this companion uses \emph{polynomial $q$}, \emph{values of
$q$}, \emph{critical values}, \emph{zero-cycle}, and \emph{moment
identities}.  The word \emph{phase} is reserved for genuine stationary-phase
or oscillatory-asymptotic arguments and is not used for the basic algebraic
data here.  The standard word \emph{constellation} is used only for the preimage
of a star joining a regular value to the critical values.  When extra arms
are added to reach $q(\operatorname{supp}\Delta)$, we refer simply to the
enlarged tree $T$ and its lifted graph $q^{-1}(T)$.

\section{Audit of the earlier draft}

The audited source is the August 2026 manuscript
\texttt{algebraic\_exponential\_integrals.tex}, containing 1,832 source
lines.  The revised base source contains about 710 lines, including
its expanded bibliography, and compiles to nine pages in the present format.

\subsection{Verdict}

No fatal gap was found in the main proof chain
\[
\begin{gathered}
\text{Beukers lifting}
\Longrightarrow
\text{constant splitting}
\Longrightarrow
\text{exponential polynomial},\\
\text{exponential polynomial}
\Longrightarrow
\text{weighted moments}
\Longrightarrow
\text{Cauchy-transform vanishing criterion}.
\end{gathered}
\]
The highest-risk points remain the standard but convention-sensitive
Fourier--Laplace identification, the localization of simple Weyl modules,
and the lateral-jump argument on a ramified lifted graph.  They are recorded
in expanded form below so that they can be audited independently.

The principal problem in the earlier draft was not a false theorem but a
mismatch between two goals: it used the foundational theorem as the vehicle
for developing a much larger structural toolkit.  The base proof was therefore
longer than necessary, while useful generalization mechanisms were buried
inside a theorem that did not require their full strength.

\subsection{Parameter absorptions and genuine simplifications}

\begin{longtable}{>{\raggedright\arraybackslash}p{0.21\textwidth}>{\raggedright\arraybackslash}p{0.34\textwidth}>{\raggedright\arraybackslash}p{0.36\textwidth}}
\toprule
Item & Earlier treatment & Streamlined treatment \\
\midrule
\endhead
Evaluation parameter $\xi$ & Carried through the main theorem and the proof. &
Replace $q$ by $\xi q$ and prove only the value at $z=1$.  Since
$I_{\xi q,\Delta}(z)=I_{q,\Delta}(\xi z)$, no information is lost. \\
\addlinespace
Affine normalization of $q$ & Translate to remove the $x^{n-1}$ term and
scale the leading coefficient to $1/n$. & Unnecessary.  For arbitrary
leading coefficient, Euclidean division still gives
$h_j=x^{j+1}/n+O(x^j)$, hence the same residue spectrum. \\
\addlinespace
Normalization of the inhomogeneous term & Define scaled vectors
$n(h_0,\ldots,h_{n-2})$ and use the prefactor $1/(nz)$. & Define
$\mathbf h=(h_j)$ directly and
write $z^{-1}\sum\mathbf b_\lambda e^{\lambda z}$.  The factor $n$ has no
mathematical role. \\
\addlinespace
Special identity $h_0=x/n$ & Used to motivate an unnecessary scaling of the inhomogeneous term. &
Not needed anywhere in the rigidity argument. \\
\addlinespace
General Beukers statement & State the full homogeneous-polynomial lifting
theorem. & Only the linear lifting statement at $z=1$ is used. \\
\addlinespace
Horizontal endomorphisms & Compute the complete algebra and deduce
multiplicity-free semisimplicity. & The base theorem needs only that a
horizontal projector is constant.  The stronger theorem is retained here. \\
\addlinespace
Cauchy geometry & Prove an exact Cauchy-transform characterization of all
functional relations. & The base theorem needs only the implication from
weighted moments to functional vanishing.  The equivalence is retained
here. \\
\addlinespace
Height bookkeeping & A separate joint-height lemma and a fully quantified
denominator construction. & The base paper uses the exact coefficient formula
plus a one-paragraph exponential denominator estimate.  The full certificate
is retained in \cref{sec:height}. \\
\bottomrule
\end{longtable}

\subsection{Corrections of emphasis and notation}

\begin{enumerate}[leftmargin=2.2em]
\item The base theorem should be advertised as an exact \emph{transfer} of
value relations to functional relations.  An explicit closed-form
classification of every functional relation is a separate geometric problem.
The distinct-regular-value criterion gives an important sufficient
case, not a general combinatorial classification.

\item The symbols used for affine scaling in the old normalization argument
collided with the support-point symbol $\alpha$.  Removing the normalization also
removes this avoidable notational hazard.

\item The semisimplicity proof is correct only after the reduced constant
summand is tracked through Fourier transform: the constant connection becomes
a module supported at the Fourier origin and vanishes after generic
localization.  This point should be stated, not left implicit.

\item Generic localization of a simple Weyl module must be justified.  Exactness
alone does not say that a nonzero localization remains simple.  The needed
one-line argument is included in \cref{prop:semisimple-expanded}.

\item The Cauchy-transform proof requires integrability at ramified vertices.
The estimate $x'(t)=O(|t-c|^{1/r-1})$ is load-bearing and should remain visible
in any compressed version.
\end{enumerate}

\section{Intrinsic moment module without normalization}

Let $q\in\Qbar[x]$ have degree $n\ge2$, put $K=\Qbar(z)$, and define
\begin{equation}\label{eq:twisted-module}
  \Mcal_q=K[x]\big/(\partial_x+zq'(x))K[x].
\end{equation}
The quotient is a $K$-vector space, not a quotient ring.  Write $[p]$ for
the class of $p$.

\begin{lemma}[Brieskorn basis]\label{lem:Brieskorn}
The classes $[1],[x],\ldots,[x^{n-2}]$ form a $K$-basis of $\Mcal_q$.
Moreover,
\[
  \C[z,z^{-1},x]\big/
  (\partial_x+zq')\C[z,z^{-1},x]
\]
is free over $\C[z,z^{-1}]$ with the same basis.  Differentiation in $z$
defines the connection
\[
  \nabla_{\partial_z}[p]=[\partial_z p+qp].
\]
\end{lemma}

\begin{proof}
The identity $[zq'p]=-[p']$ reduces every monomial of degree at least
$n-1$.  If $p\ne0$, then $(\partial_x+zq')p$ has degree
$\deg p+n-1$, so its leading term cannot cancel.  Hence the image contains
no nonzero polynomial of degree below $n-1$.  Finally,
$[\partial_z+q,\partial_x+zq']=0$, so the connection is well defined.
\end{proof}

For $0\le j\le n-2$, divide
\begin{equation}\label{eq:division-toolkit}
  qx^j=h_jq'+r_j,
  \qquad \deg r_j<n-1,
\end{equation}
and define
\[
  \mathsf C(x^j)=r_j,
  \qquad
  \mathsf D(x^j)=-h_j'.
\]
Then $\mathsf C$ is multiplication by $q$ on the Jacobian algebra
$\C[x]/(q')$, and the connection matrix is
$\mathsf C+\mathsf D/z$.

\begin{lemma}[Unnormalized residue spectrum]\label{lem:unnormalized-D}
For every polynomial $q$ of degree $n$, without any translation or scaling,
$\mathsf D$ is triangular in the monomial basis with diagonal
\[
  -\frac1n,-\frac2n,\ldots,-\frac{n-1}{n}.
\]
\end{lemma}

\begin{proof}
If $q(x)=cx^n+O(x^{n-1})$, then $q'(x)=ncx^{n-1}+O(x^{n-2})$.
Consequently,
$h_j=x^{j+1}/n+O(x^j)$ and
$-h_j'=-(j+1)x^j/n+O(x^{j-1})$.
\end{proof}

\section{Full arithmetic certificate for the moment functions}
\label{sec:height}

Fix a finite degree-zero zero-cycle
$\Delta=\sum_{\alpha\in S}d_\alpha[\alpha]$ with coefficients in $\Qbar$,
and choose a finite piecewise $C^1$ one-chain $\Gamma$ with
$\partial\Gamma=\Delta$.  Put
\[
  U_j(z)=\int_\Gamma x^j e^{zq(x)}\,dx
  =\sum_{m\ge0}u_{j,m}\frac{z^m}{m!}.
\]

\begin{lemma}[Joint height from one denominator]\label{lem:joint-height}
Let $K_0$ be a number field and let
$\alpha_0,\ldots,\alpha_m\in K_0$.  If
$\delta\alpha_k\in\mathcal O_{K_0}$ for every $k$, then
\[
 h(\alpha_0,\ldots,\alpha_m)
 \le \log\delta+
 \log\max\left(1,\max_{k,\sigma}|\sigma(\alpha_k)|\right).
\]
\end{lemma}

\begin{proof}
Use the projective coordinates
$(1:\alpha_0:\cdots:\alpha_m)=(\delta:\delta\alpha_0:\cdots:\delta\alpha_m)$.
All finite-place contributions are at most $1$, and the infinite-place
contribution is bounded by the displayed quantity.
\end{proof}

\begin{proposition}[Explicit $E$-function certificate]\label{prop:E-certificate}
Each $U_j$ is an $E$-function.  More precisely, there is a constant
$\kappa_j$ such that
\[
  h(u_{j,0},\ldots,u_{j,m})\le\kappa_j(m+1).
\]
\end{proposition}

\begin{proof}
Write
\[
  q(x)^m=\sum_{k=0}^{nm}\gamma_{m,k}x^k.
\]
Since a primitive of $x^{j+k}$ is $x^{j+k+1}/(j+k+1)$,
\begin{equation}\label{eq:exact-coefficient}
  u_{j,m}=\sum_{k=0}^{nm}\gamma_{m,k}
  \frac{\sum_{\alpha\in S}d_\alpha\alpha^{j+k+1}}{j+k+1}.
\end{equation}
Choose a number field $K_0$ containing all data and an integer $d\ge1$ such
that the coefficients of $dq$, all $d\alpha$, and all $dd_\alpha$ are
integral.  Then
\[
  \delta_m=
  d^{(n+1)m+j+2}\operatorname{lcm}(1,\ldots,nm+j+1)
\]
is a common denominator for $u_{j,0},\ldots,u_{j,m}$.  Indeed,
$d^m\gamma_{m,k}$ is integral, and the spare factor $d^{nm-k}$ makes the
exponent uniform in \eqref{eq:exact-coefficient}.

For archimedean embeddings, choose a disk containing every conjugate of every
support point.  Integrating the conjugate polynomial along the straight segment
from $0$ to the support point gives
$|\sigma(u_{j,m})|\le C^{m+j+2}$.  Since
$\log\operatorname{lcm}(1,\ldots,N)=O(N)$, \cref{lem:joint-height} gives the
height bound.

The integration-by-parts system
\[
  \mathbf U'=\left(C+\frac Dz\right)\mathbf U
  +\frac1z\sum_\lambda\mathbf b_\lambda e^{\lambda z}
\]
shows that the $\Qbar(z)$-span of the moments together with the exponentials $e^{\lambda z}$ for
$\lambda\in q(\operatorname{supp}\Delta)$ is
stable under differentiation.  Hence each $U_j$ is holonomic, completing the
$E$-function certificate.
\end{proof}

\section{Expanded Fourier--Laplace semisimplicity}

Let $\Sigma$ be the critical-value set, $V_0=\Aone_t\setminus\Sigma$, and
let $\Lcal$ be the permutation local system of the unramified cover
$q^{-1}(V_0)\to V_0$.  Let
$\Lcal^0=\ker(\operatorname{Tr}:\Lcal\to\C_{V_0})$.

\begin{proposition}[Expanded semisimplicity proof]
\label{prop:semisimple-expanded}
The differential module over $\C(z)$ with matrix
$\mathsf C+\mathsf D/z$ is semisimple.
\end{proposition}

\begin{proof}
The cover is connected, so its monodromy acts transitively on a fiber.
The permutation representation has finite image and splits as
\[
  \Lcal\simeq\C_{V_0}\oplus\Lcal^0.
\]
Maschke's theorem makes $\Lcal^0$ semisimple.  Since $q$ is finite, it is
proper and small.  Therefore
\[
  Rq_*\C_{\Aone_x}[1]\simeq\operatorname{IC}(\Lcal),
\]
and the decomposition of $\Lcal$ into irreducibles gives a decomposition of
$q_+\Ocal_{\Aone_x}$ into simple regular holonomic modules under the
Riemann--Hilbert correspondence.

Use the Fourier convention
$t\mapsto-\partial_\tau$, $\partial_t\mapsto\tau$.  Fourier transform is an
exact autoequivalence of Weyl modules and preserves simple objects.  The
constant connection transforms to a module supported at $\tau=0$ and hence
disappears after localization to $\C(\tau)$.

We record the localization point.  Let $M$ be a simple Weyl module and
$S=\C[\tau]\setminus\{0\}$.  The $S$-torsion of $M$ is a Weyl submodule:
if $p(\tau)m=0$, then $p(\tau)^2\partial_\tau m=0$.  Hence either
$S^{-1}M=0$ or $M$ injects into $S^{-1}M$.  In the latter case, every nonzero
differential submodule of $S^{-1}M$ meets $M$ after clearing a denominator;
the intersection is a nonzero Weyl submodule and therefore equals $M$.
Thus every nonzero localized simple summand remains irreducible.

To identify the localized transform, factor $q$ as its graph embedding
followed by projection.  The graph generator $\delta_q$ satisfies
\[
  (t-q(x))\delta_q=0,
  \qquad
  (\partial_x+q'(x)\partial_t)\delta_q=0.
\]
After Fourier transform it is the exponential generator $e^{-\tau q(x)}$.
On coefficient polynomials, the relative differential and connection are
\[
  \partial_x-\tau q'(x),
  \qquad
  \partial_\tau-q(x).
\]
The relative de Rham complex has no negative-degree cohomology because
$-\tau q'p$ has larger $x$-degree than $p'$ for $p\ne0$.  Its degree-zero
cohomology is
\[
  \C(\tau)[x]\big/(\partial_x-\tau q')\C(\tau)[x].
\]
Pulling back by $\tau=-z$ gives \eqref{eq:twisted-module} with connection
$\partial_z+q$.  By \cref{lem:Brieskorn}, its matrix is
$\mathsf C+\mathsf D/z$.  It is therefore a direct sum of irreducible
differential modules.
\end{proof}

\section{The complete horizontal endomorphism algebra}

For a critical point $\tau$, put
$m_\tau=\ord_\tau q'$ and $\lambda_\tau=q(\tau)$, and define
\[
  g_{\tau,k}(x)=\frac{q'(x)}{(x-\tau)^k},
  \qquad 1\le k\le m_\tau.
\]

\begin{lemma}[Critical basis and compressed residue]\label{lem:critical-full}
The $g_{\tau,k}$ form a basis of $\Vcal=\C[x]_{<n-1}$.  They diagonalize
$\mathsf C$:
\[
  \mathsf Cg_{\tau,k}=\lambda_\tau g_{\tau,k}.
\]
For a fixed critical value $\lambda$, the compression of $\mathsf D$ to the
$\lambda$-eigenspace preserves each critical-point summand.  On the
$\tau$-summand its triangular diagonal is
\[
  -\frac{m_\tau}{m_\tau+1},
  -\frac{m_\tau-1}{m_\tau+1},\ldots,
  -\frac1{m_\tau+1}.
\]
\end{lemma}

\begin{proof}
The Chinese remainder theorem gives
\[
  \C[x]/(q')\simeq
  \bigoplus_\tau\C[x]/(x-\tau)^{m_\tau}.
\]
For fixed $\tau$, the images of $g_{\tau,k}$ vanish in all other factors and
form a triangular basis in the $\tau$-factor.  Since
$q-\lambda_\tau$ vanishes to order $m_\tau+1$,
\[
  qg_{\tau,k}
  =\frac{q-\lambda_\tau}{(x-\tau)^k}q'
   +\lambda_\tau g_{\tau,k},
\]
which proves the eigenvalue identity.

Write $y=x-\tau$ and
$q-\lambda_\tau=ay^{m_\tau+1}+O(y^{m_\tau+2})$.  Then
\[
  g_{\tau,k}=a(m_\tau+1)y^{m_\tau-k}+O(y^{m_\tau-k+1}),
\]
while
\[
 -\left(\frac{q-\lambda_\tau}{(x-\tau)^k}\right)'
 =-a(m_\tau+1-k)y^{m_\tau-k}+O(y^{m_\tau-k+1}).
\]
This gives the diagonal.  If $\sigma\ne\tau$ has the same critical value,
the displayed quotient vanishes at $\sigma$ to order $m_\sigma+1$; its
derivative therefore vanishes modulo $(x-\sigma)^{m_\sigma}$, excluding
cross terms.
\end{proof}

\begin{theorem}[Horizontal-endomorphism rigidity]
\label{thm:full-endomorphism}
If $P\in M_{n-1}(\C(z))$ satisfies
\[
  P'=\left[C+\frac Dz,P\right],
\]
then $P$ is constant.  Hence the horizontal endomorphism algebra is
\[
  Z(C)\cap Z(D).
\]
The same assertion holds for the dual module.
\end{theorem}

\begin{proof}
Finite nonzero poles are excluded by pole order.  At zero, a pole of order
$m$ gives $-mP_{-m}=[D,P_{-m}]$, impossible because the spectrum of
$\operatorname{ad}D$ lies in $(-1,1)$.  Thus $P$ is polynomial.

Transpose and write $Q=Q_d z^d+Q_{d-1} z^{d-1}+\cdots$.  If $d\ge1$, the top
two coefficient equations are
\[
  [Q_d,\mathsf C]=0,
  \qquad
  dQ_d=[Q_{d-1},\mathsf C]+[Q_d,\mathsf D].
\]
Projecting to a nonzero critical-value block of $Q_d$ and using
\cref{lem:critical-full} yields
\[
  dQ_{d,\lambda}=[Q_{d,\lambda},\overline D_\lambda].
\]
Every eigenvalue of the commutator on the right is a difference of two
numbers in $(-1,0)$ and therefore lies in $(-1,1)$.  It cannot equal the
positive integer $d$.  Hence $d=0$.  Substitution gives
$[C,P]=[D,P]=0$.
\end{proof}

\begin{corollary}[Multiplicity-free semisimplicity]
The horizontal endomorphism algebra is commutative.  Consequently, the
semisimple moment connection is multiplicity-free.
\end{corollary}

\begin{proof}
The matrix $D$ has distinct eigenvalues, so its centralizer is commutative.
A repeated simple constituent of a semisimple differential module would
contribute a noncommutative full matrix algebra to the endomorphism algebra.
\end{proof}

\section{The relation submodule and its constant projector}

Let $\mathbf U$ satisfy the inhomogeneous moment system and put
\[
  \Ecal=\Span_{\C(z)}\{e^{\lambda z}:\lambda\in\Lambda\}.
\]
Define
\[
  \Rcal=
  \{\boldsymbol\phi\in\C(z)^{1\times(n-1)}:
    \boldsymbol\phi\mathbf U\in\Ecal\}.
\]

\begin{lemma}[Differential submodule of relations]
The space $\Rcal$ is a differential submodule of the dual homogeneous
connection.
\end{lemma}

\begin{proof}
Write $A=C+D/z$ and
$\mathbf f=z^{-1}\sum_\lambda\mathbf b_\lambda e^{\lambda z}$.  Then
\[
  (\boldsymbol\phi\mathbf U)'
  =(\boldsymbol\phi'+\boldsymbol\phi A)\mathbf U
   +\boldsymbol\phi\mathbf f.
\]
The left side and the last term belong to $\Ecal$ whenever
$\boldsymbol\phi\in\Rcal$.
\end{proof}

\begin{proposition}[Constant projector mechanism]
The relation module has the form
\[
  \Rcal=W\otimes_\C\C(z)
\]
for a constant subspace $W$ invariant under $C$ and $D$.
\end{proposition}

\begin{proof}
By \cref{prop:semisimple-expanded}, the dual moment module is semisimple, so
$\Rcal$ has a horizontal complement and a horizontal idempotent projector.
By \cref{thm:full-endomorphism}, the projector is constant.  Its image is the
required $W$.
\end{proof}

\begin{remark}[Why this mechanism should be preserved]
The base theorem uses the projector only once, to pass from a lifted row whose
value at $1$ is $\mathbf e_0$ to the constant row $\mathbf e_0$ itself.  In
larger systems of exponential period functions, however, the same
projector identifies the full differential submodule of functional relations.
It is therefore a natural interface with uniform relation algorithms and
with decompositions into summands detected or annihilated by the selected
period functions.
\end{remark}

\section{Exact Cauchy-transform criterion}

Let $\Delta$ be a finite degree-zero zero-cycle with complex coefficients,
let $\Gamma$ be a finite piecewise $C^1$ one-chain with
$\partial\Gamma=\Delta$, and choose a finite tree $T$ containing all critical values of $q$ and all values $q(\alpha)$ with
$\alpha\in\operatorname{supp}\Delta$.  Choose a lifted graph chain
$\widetilde\Gamma\subset q^{-1}(T)$ with the same boundary.  On an oriented
open edge $e$, write $\widetilde\Gamma$ as a linear combination of inverse
branches $x_{e,i}(t)$ and put
\[
  \psi_e(t)=\sum_i f_{e,i}x_{e,i}'(t)
  =\sum_i\frac{f_{e,i}}{q'(x_{e,i}(t))}.
\]
Thus $\psi_e(t)\,dt$ is the density contributed by the lifts of $e$.

\begin{theorem}[Exact Cauchy-transform criterion]
\label{thm:edge-exact}
For fixed $T$ and $\widetilde\Gamma$,
\[
  I_{q,\Delta}(z)\equiv0
  \quad\Longleftrightarrow\quad
  \psi_e\equiv0\text{ on every open edge }e.
\]
More generally, if $0\ne H\in\C[t]$ and
\[
  \int_\Gamma q(x)^mH(q(x))\,dx=0
  \qquad(m\ge0),
\]
then every $\psi_e$ vanishes.
\end{theorem}

\begin{proof}
For $s\notin T$, define
\[
  \mathcal C_H(s)=
  \int_{\widetilde\Gamma}\frac{H(q(x))}{s-q(x)}\,dx.
\]
The moment identities make its Laurent expansion at infinity zero.  Since
$\C\setminus T$ is connected, $\mathcal C_H$ vanishes there.  Edgewise
change of variables gives
\[
  \mathcal C_H(s)=
  \sum_e\int_e\frac{H(t)\psi_e(t)}{s-t}\,dt.
\]
At an interior point of one edge, contributions from all other edges extend
holomorphically across a small disk.  The local Cauchy integral has lateral
jump
$2\pi i\varepsilon_e H(t)\psi_e(t)$.  Both lateral values vanish, so
$H(t)\psi_e(t)=0$.  Away from the finite zero set of $H$, this gives
$\psi_e=0$, and analyticity extends the conclusion to the entire edge.
The exponential integral is the edgewise Laplace transform of the densities,
proving one direction.  The other is immediate.  Taking $H=1$ after
$I_{q,\Delta}\equiv0$ gives the converse implication in the equivalence.
This is the standard Cauchy-transform mechanism used in the polynomial
moment literature; compare \cite{GavrilovPakovich,PakovichRoytvarfYomdin}.
\end{proof}

\begin{corollary}[Functional independence over distinct regular values]
If the support points $a_0,\ldots,a_r$ have pairwise distinct regular
values under $q$, then the zero-cycle
$\sum_i d_i[a_i]$ has identically zero exponential integral only when all
$d_i=0$.
\end{corollary}

\begin{proof}
Give each $q(a_i)$ a separate terminal arm.  On that arm, the only nonzero
terminal coefficient is $d_i$, and the density equals
$d_i/q'(x_i(t))$.  Apply \cref{thm:edge-exact}.
\end{proof}

\section{What is likely to generalize}

The following division should guide later papers on exponential
period functions.

\subsection{Robust structural mechanisms}

\begin{enumerate}[leftmargin=2.2em]
\item \textbf{Differential submodule of relations and projector.}
This is not tied to a scalar integral.  The same construction applies
to an inhomogeneous differential system whose homogeneous module is
semisimple and whose inhomogeneous term lies in a
finite-dimensional differential submodule generated by exponential functions.

\item \textbf{Nonresonant residue criterion.}
The pole arguments use only that residue eigenvalue differences avoid
nonzero integers, locally at finite points and after compression at infinity.
This suggests an abstract nonresonant constant-splitting criterion.

\item \textbf{Fourier--Laplace semisimplicity.}
Finite monodromy before Fourier transform is a particularly transparent
source of semisimplicity.  More general semisimple local systems on curves
should produce analogous transformed modules.

\item \textbf{Cauchy jump detection.}
Moment vanishing is converted into vanishing of a Cauchy transform, and local
jumps recover the density.  The argument should extend to algebraic finite
covers of curves once a suitable cut graph and local branches are available.
\end{enumerate}

\subsection{Polynomial-specific pieces}

\begin{enumerate}[leftmargin=2.2em]
\item The monomial Brieskorn basis $1,x,\ldots,x^{n-2}$ and the division
algorithm by $q'$ use the affine line essentially.

\item The critical basis $q'/(x-\tau)^k$ and the explicit fractions
$-r/(m+1)$ are special to one-variable polynomial ramification.

\item The chosen tree $T$ lies in the target plane.  On a positive-genus
curve it must be replaced by a cut graph, and closed-cycle periods may remain
after all edge densities vanish.
\end{enumerate}

\subsection{Algorithmic outputs worth preserving}

For a concrete polynomial $q$, the following data are explicit and small:
\[
  C,\quad D,\quad
  \{(\tau,m_\tau,\lambda_\tau)\},\quad
  \{\overline D_\lambda\},\quad
  \{\mathbf b_\lambda\}.
\]
They provide certificates for residue nonresonance, constant projectors, and
detection of boundary coefficients without reconstructing a large minimal
scalar operator.  This is the natural point of contact with the
Adamczewski--Rivoal and Bostan--Rivoal--Salvy exceptional-value algorithms.

\section*{AI-assisted research disclosure}

This companion was developed through interactive work between Christopher
D. Long and ChatGPT 5.6 Pro.  ChatGPT assisted in proof compression,
line-by-line auditing, symbolic stress tests, organization, and drafting.
The AI system is not an author.  The human author bears full responsibility
for all mathematical claims and any remaining errors.

\small
\begin{thebibliography}{99}

\bibitem{AdamczewskiRivoal}
B.~Adamczewski and T.~Rivoal,
\emph{Exceptional values of $E$-functions at algebraic points},
Bull. Lond. Math. Soc. \textbf{50} (2018), no.~4, 697--708.
\href{https://doi.org/10.1112/blms.12168}{doi:10.1112/blms.12168}.

\bibitem{Beukers}
F.~Beukers,
\emph{A refined version of the Siegel--Shidlovskii theorem},
Ann. of Math. (2) \textbf{163} (2006), no.~1, 369--379.
\href{https://doi.org/10.4007/annals.2006.163.369}{doi:10.4007/annals.2006.163.369}.

\bibitem{BostanRivoalSalvy}
A.~Bostan, T.~Rivoal, and B.~Salvy,
\emph{Minimization of differential equations and algebraic values of
$E$-functions},
Math. Comp. \textbf{93} (2024), no.~347, 1427--1472.
\href{https://doi.org/10.1090/mcom/3912}{doi:10.1090/mcom/3912}.

\bibitem{CommelinHabeggerHuber}
J.~Commelin, P.~Habegger, and A.~Huber,
\emph{Exponential periods and o-minimality},
arXiv:2007.08280, version of August 11, 2026.

\bibitem{deCataldoMigliorini}
M.~A.~A. de~Cataldo and L.~Migliorini,
\emph{The decomposition theorem, perverse sheaves and the topology of
algebraic maps},
Bull. Amer. Math. Soc. (N.S.) \textbf{46} (2009), no.~4, 535--633.
\href{https://doi.org/10.1090/S0273-0979-09-01260-9}{doi:10.1090/S0273-0979-09-01260-9}.

\bibitem{FresanJossen}
J.~Fres\'an and P.~Jossen,
\emph{A non-hypergeometric $E$-function},
Ann. of Math. (2) \textbf{194} (2021), no.~3, 903--942.
\href{https://doi.org/10.4007/annals.2021.194.3.7}{doi:10.4007/annals.2021.194.3.7}.

\bibitem{GavrilovPakovich}
L.~Gavrilov and F.~Pakovich,
\emph{Moments on Riemann surfaces and hyperelliptic Abelian integrals},
Comment. Math. Helv. \textbf{89} (2014), no.~1, 125--155.
\href{https://doi.org/10.4171/CMH/314}{doi:10.4171/CMH/314}.

\bibitem{HienRoucairol}
M.~Hien and C.~Roucairol,
\emph{Integral representations for solutions of exponential Gauss--Manin
systems},
Bull. Soc. Math. France \textbf{136} (2008), no.~4, 505--532.
\href{https://doi.org/10.24033/bsmf.2564}{doi:10.24033/bsmf.2564}.

\bibitem{HTT}
R.~Hotta, K.~Takeuchi, and T.~Tanisaki,
\emph{$D$-Modules, Perverse Sheaves, and Representation Theory},
Progress in Mathematics, vol.~236,
Birkh\"auser Boston, Boston, MA, 2008.
\href{https://doi.org/10.1007/978-0-8176-4523-6}{doi:10.1007/978-0-8176-4523-6}.

\bibitem{PakovichMuzychuk}
F.~Pakovich and M.~Muzychuk,
\emph{Solution of the polynomial moment problem},
Proc. Lond. Math. Soc. (3) \textbf{99} (2009), no.~3, 633--657.
\href{https://doi.org/10.1112/plms/pdp010}{doi:10.1112/plms/pdp010}.

\bibitem{PakovichRoytvarfYomdin}
F.~Pakovich, N.~Roytvarf, and Y.~Yomdin,
\emph{Cauchy-type integrals of algebraic functions},
Israel J. Math. \textbf{144} (2004), 221--291.
\href{https://doi.org/10.1007/BF02916714}{doi:10.1007/BF02916714}.

\bibitem{SabbahYu}
C.~Sabbah and J.-D.~Yu,
\emph{On the irregular Hodge filtration of exponentially twisted mixed
Hodge modules},
Forum Math. Sigma \textbf{3} (2015), Paper No.~e9, 71 pp.
\href{https://doi.org/10.1017/fms.2015.5}{doi:10.1017/fms.2015.5}.

\end{thebibliography}
\end{document}
